\section{Dos atolladeros: flotar indefinidamente y que las cuentas no cuadren}

La sección anterior trató el software, el hardware y los datos; esta sección aborda los riesgos de las empresas emergentes. Wang He (王鹤) sostiene que, si una empresa de inteligencia corporizada cae en alguno de dos atolladeros, su techo será muy limitado. El primero es «flotar indefinidamente»: la empresa permanece mucho tiempo en el terreno de los conceptos, las demostraciones y el relato para atraer financiamiento, sin entrar en escenarios reales ni someterse a la prueba de los clientes, los costos y la confiabilidad. El segundo es que «las cuentas no cuadren»: aun cuando entra en un escenario, el costo marginal no baja, cada entrega pesa más que la anterior y, al final, la productividad no se materializa.

En esencia, ambos atolladeros son un desajuste en la retroalimentación. Flotar indefinidamente significa que el reward model de la empresa es la attention y la valuation; que las cuentas no cuadren significa que la ruta tecnológica de la empresa no ha logrado reducir el costo unitario de entrega. La inteligencia corporizada es una industria física: no basta con mirar el benchmark del modelo ni los videos demo. Tiene que responder lo siguiente: una vez que un robot entra en un escenario, ¿puede completar un tipo de tarea de manera más estable, más barata o más escalable que una persona?

\begin{figure}[H]
\centering
\includegraphics[width=0.94\textwidth]{two-traps.png}
\caption{Dos atolladeros: tanto flotar indefinidamente como que las cuentas no cuadren limitan el techo de la empresa. Diagrama conceptual de elaboración propia, a partir de la conversación entre 01:44:34 y 01:55:17 y de la descripción del video.}
\end{figure}

\begin{knowledgebox}{Lectura de la figura: uno es un riesgo de relato y el otro, un riesgo económico}
El problema de flotar indefinidamente es desconectarse de los clientes reales; el problema de que las cuentas no cuadren es que cuanto más se entrega, más se pierde. Al leer la figura hay que relacionar ambos lados: hacer solo demos evita exponer los costos a corto plazo, pero también hace perder los datos reales; y si, una vez dentro del escenario, el costo marginal no baja, el volante de datos tampoco puede sostener el ciclo comercial cerrado.
\end{knowledgebox}

\subsection{Generalización dentro del escenario de aplicación}

Esta sección explica por qué Galbot, Yinhe Tongyong (银河通用), eligió la «generalización dentro del escenario de aplicación». Los robots de propósito general resultan atractivos, pero buscar demasiado pronto la generalización a todos los escenarios hace que las tareas, el hardware, los datos y los costos se disparen al mismo tiempo. Una ruta más realista es generalizar a fondo dentro de un escenario de aplicación concreto: por ejemplo, en el escenario de anaqueles, los distintos productos, su acomodo, las oclusiones, los faltantes, el reabastecimiento y las situaciones anómalas. No se trata de automatizar un solo punto, sino de construir capacidades escalables dentro de un mundo acotado.

Este enfoque se parece también a la conversión en producto de los modelos grandes: primero se logra que funcionen de punta a punta el dominio de tareas, la retroalimentación, los costos y el ciclo cerrado de datos, y luego se amplían gradualmente los límites de la capacidad. Para los robots, la generalización dentro del escenario tiene una ventaja adicional: una vez que hay unidades entregadas, los datos reales regresan y el sistema puede aprender de las fallas reales. Sin volumen de entregas, el volante de datos es solo una ilusión.

\begin{importantbox}{Generalización dentro del escenario}
En la etapa temprana de la inteligencia corporizada, la ruta eficaz quizá no sea «hacer primero un robot de propósito general», sino lograr una generalización suficiente de objetos, tareas, anomalías y costos dentro de un escenario real por el que se paga. Cuanto más claros sean los límites del escenario, más fácil será validar el ciclo cerrado de datos y el ciclo comercial cerrado.
\end{importantbox}

\subsection{Resumen de la sección}

Las empresas de inteligencia corporizada deben evitar dos atolladeros: no quedarse flotando indefinidamente en el relato y no entrar en un escenario donde las cuentas no cuadren. La ruta correcta es encontrar un escenario por el que se pague, generalizar dentro de él y lograr que las unidades entregadas devuelvan datos.
